\documentclass[11pt]{article}
\usepackage[T1]{fontenc}
\usepackage[utf8]{inputenc}
\usepackage[margin=1in]{geometry}
\usepackage{amsmath,amssymb,amsthm}
\usepackage{microtype}
\usepackage[hidelinks]{hyperref}
\newcommand{\E}{\mathbb E}
\newcommand{\Pb}{\mathbb P}
\newcommand{\Ind}[1]{\mathbf 1\{#1\}}
\newtheorem{proposition}{Proposition}
\newtheorem{lemma}{Lemma}
\title{Fluid model forgetting case}
\author{Standalone theoretical note}
\date{}
\begin{document}
\maketitle

This note presents a proposition where we can characterize the region where there is at least one threshold if the set is non empty.

Recall from the previous note 

\paragraph{Model and notation.}
Fix $\rho\in(0,1)$, $\gamma\in(0,1)$, $p_0\in(0,1)$,
$0<p_1<p_2<1$, and $0\leq c_1<c_2<R$.
Write $\Delta c=c_2-c_1$ and $\Delta p=p_2-p_1$.
As in the manuscript, the normalized state space and outcome transitions are
\begin{align*}
\mathcal Z&=\{(x_s,x_f):x_s\geq0,\ x_f\geq0,\ x_s+x_f\leq1\},\\
T_\rho^s(x_s,x_f)&=(\rho x_s+1-\rho,\rho x_f),\\
T_\rho^f(x_s,x_f)&=(\rho x_s,\rho x_f+1-\rho).
\end{align*}


\begin{proposition}[Characterization of positive advantage]
\label{prop:positive-characterization}
Let $U^{1}_\rho$ be the value of always choosing product~1.
For $\rho\in(0,1)$, we have $\mathcal A_\rho>0$ if and only if
\begin{equation}
\max_{z\in\mathcal K_\rho}
\left[U_\rho^1(T_\rho^s z)-U_\rho^1(T_\rho^f z)\right]
>\frac{\Delta c}{\gamma\Delta p}.
\end{equation}
\end{proposition}

\begin{proof}
Conditional on a purchase from seller~$A$, let
\begin{equation*}
Q_\rho^i(z)
=R-c_i+\gamma\left[p_iU_\rho^1(T_\rho^s z)
+(1-p_i)U_\rho^1(T_\rho^f z)\right],
\qquad i\in\{1,2\}.
\end{equation*}
This is the value of choosing product~$i$ for the current purchase and
product~1 thereafter. Subtracting the two expressions gives
\begin{equation*}
Q_\rho^2(z)-Q_\rho^1(z)
=-\Delta c+\gamma\Delta p
\left[U_\rho^1(T_\rho^s z)-U_\rho^1(T_\rho^f z)\right].
\end{equation*}
Define
\begin{equation*}
\mathcal A_\rho^1
=\max_{z\in\mathcal K_\rho}
\left\{-\Delta c+\gamma\Delta p
\left[U_\rho^1(T_\rho^s z)-U_\rho^1(T_\rho^f z)\right]\right\}.
\end{equation*}
The fixed-policy operator is a contraction preserving continuity,
so $U_\rho^1$ is continuous and this maximum exists.
We first show that $\mathcal A_\rho>0$ if and only if
$\mathcal A_\rho^1>0$, by proving that their nonpositivity conditions
are equivalent.

Suppose $\mathcal A_\rho^1\leq0$. Then
$ Q_\rho^1(z)\geq Q_\rho^2(z)$ at every
$z\in\mathcal K_\rho$. By the fixed-policy Bellman equation,
\begin{equation*}
U_\rho^1(z)
=\frac{D_\rho(z)}{1-\gamma[1-D_\rho(z)]}Q_\rho^1(z)
=\frac{D_\rho(z)}{1-\gamma[1-D_\rho(z)]}
\max_{i\in\{1,2\}} Q_\rho^i(z).
\end{equation*}
Thus $U_\rho^1$ also satisfies the optimal Bellman equation on
$\mathcal K_\rho$. Since this set is invariant under both transitions,
the Bellman operator restricted to it is a contraction. Uniqueness
therefore gives $U_\rho^1=V_\rho$ on $\mathcal K_\rho$, and hence
$\mathcal A_\rho=\mathcal A_\rho^1\leq0$.

Conversely, suppose $\mathcal A_\rho\leq0$. Product~1 then maximizes
the Bellman expression evaluated at $V_\rho$ at every state in
$\mathcal K_\rho$. Consequently, $V_\rho$ satisfies the Bellman equation
for always choosing product~1 on this invariant set. Uniqueness of
the fixed-policy value gives $V_\rho=U_\rho^1$ there, and hence
$\mathcal A_\rho^1=\mathcal A_\rho\leq0$.

We have proved $\mathcal A_\rho\leq0$ if and only if
$\mathcal A_\rho^1\leq0$. Taking negations gives the desired equivalence
of strict positivity. Finally, since $\Delta c$ and $\gamma\Delta p>0$
do not depend on $z$,
\begin{equation*}
\mathcal A_\rho^1
=-\Delta c+\gamma\Delta p\max_{z\in\mathcal K_\rho}
\left[U_\rho^1(T_\rho^s z)-U_\rho^1(T_\rho^f z)\right].
\end{equation*}
Therefore $\mathcal A_\rho^1>0$ is equivalent to the inequality
in the statement.
\end{proof}

Now we can think about the information that the fix policy of using product 2 gives us. Only because it is easier to read let's have the following notation
\begin{align}
    a_\rho(z) &= -\Delta c + \gamma \Delta p \left[V_\rho(T_\rho^s z)-V_\rho(T_\rho^f z)\right] \\
    a^1_\rho(z) &= -\Delta c + \gamma \Delta p\left[U_\rho^1(T_\rho^s z)-U_\rho^1(T_\rho^f z)\right] \\
    a^2_\rho(z) &= -\Delta c + \gamma \Delta p\left[U_\rho^2(T_\rho^s z)-U_\rho^2(T_\rho^f z)\right]
\end{align}

Then $\mathcal A_\rho = \max_{z\in\mathcal K_\rho} a_\rho(z)$ and $\mathcal A_\rho^1 = \max_{z\in\mathcal K_\rho} a^1 _\rho(z)$. We define 
\begin{equation}
    \mathcal A_\rho^2 = \min_{z\in\mathcal K_\rho} a^2_\rho(z)
\end{equation}

We can characterize the region where there exist at least one threshold thinking about the intersection of states where there is used of product 2 based on conditions determined over $\mathcal A_\rho^1$ and $\mathcal A_\rho^2$. 

\begin{proposition}
\label{prop:interior-indifference}
For $\rho\in(\frac12,1)$, if $\mathcal A_\rho^1>0$ and
$\mathcal A_\rho^2<0$, then there exists $z^*\in\mathcal K_\rho$
such that $a_\rho(z^*)=0$.
\end{proposition}

\begin{proof}
By proposition \ref{prop:positive-characterization} we have that $\mathcal A_\rho^1>0$ implies that there exist $z_1 \in \mathcal K_\rho$ such that $a_\rho(z_1)>0$. 
We can prove that $\mathcal A_\rho^2<0 \iff \min_{z\in\mathcal K_\rho} a_\rho(z) <0 $. Indeed, suppose that $\mathcal A_\rho^2\ge0$ then, for all $z \in \mathcal K_\rho$
\begin{equation}
    -\Delta c + \gamma \Delta p\left[U_\rho^2(T_\rho^s z)-U_\rho^2(T_\rho^f z)\right] \ge 0 
\end{equation}
Hence, product~2 maximizes the Bellman expression evaluated at
$U_\rho^2$ at every state in $\mathcal K_\rho$. By the fixed-policy
Bellman equation,
\begin{align*}
U_\rho^2(z)
&=\frac{D_\rho(z)}{1-\gamma[1-D_\rho(z)]}
\left\{R-c_2+\gamma\left[p_2U_\rho^2(T_\rho^s z)
+(1-p_2)U_\rho^2(T_\rho^f z)\right]\right\}\\
&=\frac{D_\rho(z)}{1-\gamma[1-D_\rho(z)]}
\max_{i\in\{1,2\}}
\left\{R-c_i+\gamma\left[p_iU_\rho^2(T_\rho^s z)
+(1-p_i)U_\rho^2(T_\rho^f z)\right]\right\}.
\end{align*}
Thus $U_\rho^2$ also satisfies the optimal Bellman equation on
$\mathcal K_\rho$. Since this set is invariant under both transitions,
the Bellman operator restricted to it is a contraction. Uniqueness
therefore gives $U_\rho^2=V_\rho$ on $\mathcal K_\rho$.
Both successor states remain in this set, so $a_\rho^2=a_\rho$ there,
and hence
\begin{equation*}
\min_{z\in\mathcal K_\rho}a_\rho(z)
=\min_{z\in\mathcal K_\rho}a_\rho^2(z)
=\mathcal A_\rho^2\geq0.
\end{equation*}

Conversely, suppose $\min_{z\in\mathcal K_\rho}a_\rho(z)\geq0$.
Then $a_\rho(z)\geq0$ at every state in $\mathcal K_\rho$, so
product~2 maximizes the Bellman expression evaluated at $V_\rho$.
Consequently,
\begin{equation*}
V_\rho(z)
=\frac{D_\rho(z)}{1-\gamma[1-D_\rho(z)]}
\left\{R-c_2+\gamma\left[p_2V_\rho(T_\rho^s z)
+(1-p_2)V_\rho(T_\rho^f z)\right]\right\}.
\end{equation*}
Thus $V_\rho$ satisfies the Bellman equation for always choosing
product~2 on this invariant set. Uniqueness of the fixed-policy value
gives $V_\rho=U_\rho^2$ there. Therefore $a_\rho=a_\rho^2$ on
$\mathcal K_\rho$, and
\begin{equation*}
\mathcal A_\rho^2
=\min_{z\in\mathcal K_\rho}a_\rho^2(z)
=\min_{z\in\mathcal K_\rho}a_\rho(z)\geq0.
\end{equation*}

We have proved that $\mathcal A_\rho^2\geq0$ if and only if
$\min_{z\in\mathcal K_\rho}a_\rho(z)\geq0$. Taking negations gives
\begin{equation*}
\mathcal A_\rho^2<0
\quad\Longleftrightarrow\quad
\min_{z\in\mathcal K_\rho}a_\rho(z)<0.
\end{equation*}
Since $\mathcal A_\rho^2<0$, there exists $z_2\in\mathcal K_\rho$
such that $a_\rho(z_2)<0$. Together with $a_\rho(z_1)>0$, this gives
two states at which $a_\rho$ has opposite signs.

Since $\rho>1/2$, the set $\mathcal K_\rho$ is the segment
$\{(x_s,1-x_s):x_s\in[0,1]\}$. We can therefore write
$z_1=(x_1,1-x_1)$ and $z_2=(x_2,1-x_2)$, with $x_1,x_2\in[0,1]$.
Then
\begin{equation*}
a_\rho(x_2,1-x_2)<0<a_\rho(x_1,1-x_1).
\end{equation*}
Since $x\mapsto a_\rho(x,1-x)$ is continuous on $[0,1]$,
the intermediate value theorem therefore gives
\begin{equation*}
x^*\in\bigl(\min\{x_1,x_2\},\max\{x_1,x_2\}\bigr)\subseteq(0,1)
\end{equation*}
such that
\begin{equation*}
a_\rho(x^*,1-x^*)=0.
\end{equation*}
Thus $z^*=(x^*,1-x^*)\in\mathcal K_\rho$ satisfies $a_\rho(z^*)=0$.
At this state, both products give the same Bellman value and are
therefore optimal.
\end{proof}


\paragraph{A cost interval for an interior indifference state.}
For $i\in\{1,2\}$, let
\begin{equation}
H_\rho^i(z)
=\E_z^i\left[\sum_{t=1}^{\infty}\gamma^{t-1}\Ind{y_t=A}\right]
\end{equation}
be the expected discounted number of purchases from seller~$A$ when
product~$i$ is always used. Here $t$ counts calendar periods, and
$\E_z^i$ denotes expectation starting from $z$ under this fixed policy.

\begin{proposition}[Cost region with both products strictly optimal at some states]
\label{prop:cost-region}
Fix $\rho\in(1/2,1)$, $\gamma\in(0,1)$, $p_0\in(0,1)$,
$0<p_1<p_2<1$, and $0\leq c_1<R$. Define
\begin{align*}
\underline c_2
&=\frac{
 c_1+\gamma\Delta p\,R
 \displaystyle\min_{z\in\mathcal K_\rho}
 \left[H_\rho^2(T_\rho^s z)-H_\rho^2(T_\rho^f z)\right]
}{
 1+\gamma\Delta p
 \displaystyle\min_{z\in\mathcal K_\rho}
 \left[H_\rho^2(T_\rho^s z)-H_\rho^2(T_\rho^f z)\right]
},\\[1ex]
\overline c_2
&=c_1+\gamma\Delta p\,(R-c_1)
\max_{z\in\mathcal K_\rho}
\left[H_\rho^1(T_\rho^s z)-H_\rho^1(T_\rho^f z)\right].
\end{align*}
These bounds do not depend on $c_2$. For any $c_2\in(c_1,R)$,
\begin{equation}
\mathcal A_\rho^1>0\ \text{and}\ \mathcal A_\rho^2<0
\quad\Longleftrightarrow\quad
\underline c_2<c_2<\overline c_2.
\end{equation}
Thus $(\underline c_2,\overline c_2)\cap(c_1,R)$ is exactly the cost
region where each product is strictly optimal at some state in
$\mathcal K_\rho$. For every cost in this region, there exists
$x^*\in(0,1)$ such that $a_\rho(x^*,1-x^*)=0$.
\end{proposition}

\begin{proof}
Under the fixed policy of always choosing product~$i$, the distribution
of purchases and outcomes does not depend on costs. Moreover,
\begin{equation*}
0\leq H_\rho^i(z)
\leq\sum_{t=1}^{\infty}\gamma^{t-1}
=\frac{1}{1-\gamma}<\infty.
\end{equation*}
Since each purchase gives the constant profit $R-c_i$, we have
\begin{equation*}
U_\rho^i(z)=(R-c_i)H_\rho^i(z),\qquad i\in\{1,2\}.
\end{equation*}
In particular, both $H_\rho^1$ and $H_\rho^2$ are independent of $c_2$.

We first check that the denominator defining $\underline c_2$ is
positive. Conditioning on the first calendar period gives
\begin{equation*}
H_\rho^i(z)
=\frac{D_\rho(z)}{1-\gamma[1-D_\rho(z)]}
\left\{1+\gamma\left[p_iH_\rho^i(T_\rho^s z)
+(1-p_i)H_\rho^i(T_\rho^f z)\right]\right\}.
\end{equation*}
On the boundary, demand $D_\rho(x,1-x)$ and both successor coordinates
are increasing in $x$. Iteration of this fixed-policy equation from
zero therefore shows that $H_\rho^i(x,1-x)$ is nondecreasing.
It follows that
\begin{equation*}
H_\rho^2(T_\rho^s z)-H_\rho^2(T_\rho^f z)\geq0,
\qquad z\in\mathcal K_\rho,
\end{equation*}
so the denominator defining $\underline c_2$ is at least one.
The same fixed-policy operator is a contraction preserving
continuity, so the extrema in the definitions exist.

Using $\Delta c=c_2-c_1$ and the factorization of $U_\rho^1$, we obtain
\begin{align*}
\mathcal A_\rho^1
&=c_1-c_2+\gamma\Delta p\,(R-c_1)
\max_{z\in\mathcal K_\rho}
\left[H_\rho^1(T_\rho^s z)-H_\rho^1(T_\rho^f z)\right]\\
&=\overline c_2-c_2.
\end{align*}
Hence $\mathcal A_\rho^1>0$ if and only if $c_2<\overline c_2$.

Similarly, the factorization of $U_\rho^2$ gives
\begin{equation*}
\mathcal A_\rho^2
=c_1-c_2+\gamma\Delta p\,(R-c_2)
\min_{z\in\mathcal K_\rho}
\left[H_\rho^2(T_\rho^s z)-H_\rho^2(T_\rho^f z)\right].
\end{equation*}
Collecting the terms containing $c_2$, the condition
$\mathcal A_\rho^2<0$ is equivalent to
\begin{align*}
&c_1+\gamma\Delta p\,R
\min_{z\in\mathcal K_\rho}
\left[H_\rho^2(T_\rho^s z)-H_\rho^2(T_\rho^f z)\right]\\
&\qquad<c_2\left\{1+\gamma\Delta p
\min_{z\in\mathcal K_\rho}
\left[H_\rho^2(T_\rho^s z)-H_\rho^2(T_\rho^f z)\right]\right\}.
\end{align*}
Dividing by the positive expression in braces gives
$\mathcal A_\rho^2<0$ if and only if $c_2>\underline c_2$.
Combining the two inequalities proves the stated equivalence.

By Proposition~\ref{prop:positive-characterization} and the equivalence
proved in Proposition~\ref{prop:interior-indifference}, these conditions
hold exactly when $a_\rho$ is positive at some state and negative at
another. Proposition~\ref{prop:interior-indifference} then gives the
interior indifference state.
\end{proof}

\paragraph{Remark.}
The remaining parameters are arbitrary within the assumptions, but they
determine the two cost bounds. The admissible cost region is their open
interval intersected with $(c_1,R)$, whenever this intersection is
nonempty. The proposition establishes the existence of an interior
indifference state; it does not determine the number or ordering of
policy switches.

\end{document}